\documentclass{article}
\usepackage[utf8]{inputenc}
\usepackage{url}
\usepackage{graphicx}
\usepackage{blindtext}
\usepackage{hyperref}
\usepackage[margin=3cm]{geometry}
\usepackage{color}
\usepackage{tikz}
\usepackage{titlesec}
\usepackage{titling}
\usepackage{placeins}

\newcommand{\logo}{\includegraphics[width=5cm]{nurus_logo.png}}

\title{\Huge\textbf{Sponsorship Report}\\[1cm]\LARGE Project Results for NURUS}

\author{\Large\textbf{Presented by:}\\[0.5cm] Muhammed Sezer \& Şevval Belkıs Dikkaya}

\begin{document}
\begin{titlepage}
    \centering
    
    % NURUS logo
    \logo
    
    \vspace{1cm}
    
    % Title
    \maketitle
    
    \vfill
    
    % Footer
    \textbf{Contact Information:}\\[0.2cm]
    Muhammed Sezer\\
    Phone: 530-613-61-84\\
    Email: muhammedsezer12@gmail.com
\end{titlepage}

\section{Introduction}

In this report, we will detail our sponsorship agreement with NURUS, including the terms of the agreement, the benefits we received, and the results of our collaboration.
\section{Our Team}

\FloatBarrier

\begin{table}[hbt!]
\centering
\begin{tabular}{|c|c|c|}
    \hline
    Name & Major & LinkedIn \\
    \hline
    Metehan İçöz & Mechanical Engineering & \href{https://www.linkedin.com/in/metehan-içöz-735bab206/}{\includegraphics[height=1.2em]{linkedin.png}} \\
    Muhammed Sezer & Electrical - Electronics Engineering & \href{https://www.linkedin.com/in/muhammed-sezer-160428208/}{\includegraphics[height=1.2em]{linkedin.png}} \\
    Şevval Belkıs Dikkaya & Electrical - Electronics Engineering & \href{https://www.linkedin.com/in/sbdikkaya/}{\includegraphics[height=1.2em]{linkedin.png}} \\
    Ecem Ataman & Aerospace Engineering & \href{https://www.linkedin.com/in/ecemataman/}{\includegraphics[height=1.2em]{linkedin.png}} \\
    Zeynep Keleş & Electrical - Electronics Engineering & \href{https://www.linkedin.com/in/zeynep-bilge-keleş-7ba917255/}{\includegraphics[height=1.2em]{linkedin.png}} \\
    \hline
\end{tabular}
\caption[short]{Team Participants}
\end{table}


Our team has a strong background in AI, coding, design, and robotics, which allowed us to create a high-quality UAV for the Teknofest competition. As a result of our hard work and dedication, we were able to secure first and second place in the AI in Transportation competition. We are proud to have been sponsored by NURUS for this competition, as well as for the Fighter UAV competition. This sponsorship allowed us to use the highest-quality materials and equipment to create our UAV.

At our core, we are a team that values sharing and encouraging others. We believe in the power of collaboration and helping each other succeed. It is this philosophy that has driven our success not only in the competition but also in our academic and professional careers.
\FloatBarrier

\section{Background}

Our university team was tasked with designing and building a fixed-wing aircraft for Teknofest. In order to create a high-quality UAV, we needed access to specialized equipment and materials. We reached out to NURUS, an innovative and R\&D-oriented company that produces office furniture, as well as works in the fields of aviation and technology, to explore the possibility of sponsorship.

After contacting Renan Gökyay, a representative of NURUS, we were able to secure a sponsorship package that included a budget of 20.000TL and access to their workspace and equipment. We used a variety of top-quality materials, including fabric, carbon, epoxy, and other composites, to build our fixed-wing aircraft.

\section{Sponsorship Agreement}

The sponsorship agreement began in December 2022 and will end in April 2023. As part of the agreement, NURUS required us to place their logos on the aircraft and team t-shirt, and to create a video in their factory.

\section{Benefits}

Through our collaboration with NURUS, we were able to benefit from their expertise and resources, including their workspace and equipment. We also received financial support.
Thanks to NURUS' sponsorship, we had access to a wealth of opportunities within their factory, which allowed us to create a superior aircraft compared to our competitors.

\section{Factory Experience}

During our time at NURUS, we had access to their workspace and equipment, which allowed us to work on our aircraft in a professional setting. We were also able to collaborate with their team members who provided us with valuable insights and guidance on how to improve our aircraft.

Here are some of the steps we took during the process:

\begin{itemize}
    \item We used NURUS' specialized software to design and test our aircraft.
    \item We used their top-quality materials, including fabric, carbon, epoxy, and other composites, to build our aircraft.
    \item We used their tools and equipment to manufacture and assemble the aircraft, including a CNC machine, a laser cutter, and a 3D printer.
    \item We worked closely with their team members to troubleshoot any issues and optimize our design.  
	\item Access to NURUS' network of industry professionals and potential collaborators
	\item Exposure to new technologies and manufacturing processes, including additive manufacturing and composites fabrication
	\item Training sessions and workshops on topics such as design software, machining, and testing techniques
	\item Feedback and guidance from NURUS' experienced engineers and designers
	\item Opportunities to showcase your work and promote your team through NURUS' social media and marketing channels
	\item Invitations to attend industry events and conferences as representatives of NURUS and your university
	\item Potential for ongoing collaboration with NURUS on future projects or research initiatives
\end{itemize}

\section{Our Experience with NURUS}

Working with NURUS was an unforgettable experience that challenged us to push the limits of our knowledge and skills. From the moment we stepped into their workspace, we were surrounded by a team of passionate engineers and designers who were eager to help us succeed.

Our journey began with the design phase, where we used NURUS' specialized software to create a 3D model of our aircraft. With the help of their team members, we were able to refine our design and optimize it for performance and efficiency.

Once we had our design locked in, we began the manufacturing process. This was where NURUS truly shone, as we were given access to their state-of-the-art equipment and materials. We used their CNC machine to precisely cut our components, their laser cutter to create intricate patterns in our wings, and their 3D printer to produce small parts with extreme accuracy. We also used their high-quality materials, including carbon fiber, epoxy, and composites, to construct our aircraft.

As we worked on our aircraft, we were constantly supported by NURUS' team members, who provided us with invaluable feedback and guidance. We were encouraged to ask questions and seek help whenever we needed it, and we always felt like we were part of the NURUS family.

One of the most memorable moments of our experience was when we were invited to present our progress to NURUS' management team. We gave a comprehensive presentation on our design, manufacturing, and testing processes, and were met with enthusiastic applause and encouragement.

Throughout our time at NURUS, we also had the opportunity to learn about the company's other projects and initiatives, which ranged from cutting-edge furniture design to advanced aviation research. We were inspired by their commitment to innovation and excellence, and felt privileged to be part of their community.

As our project came to a close, we were proud to see our aircraft take shape and perform at its best. We knew that without NURUS' support, we would not have been able to achieve such a high level of success. We are grateful to NURUS for giving us this opportunity, and we look forward to maintaining our relationship with them in the future.

\section{Results}

As a result of our partnership with NURUS, we were able to create an exceptional aircraft that outperformed our competitors. We were also able to fulfill the requirements of the sponsorship agreement by placing NURUS' logos on our aircraft and t-shirt, and by creating a video in their factory.


\section{Conclusion}

Our sponsorship agreement with NURUS was a success, and we are grateful for the support and resources they provided us. We look forward to the possibility of future collaborations with NURUS.

\section{Media Coverage and Public Relations}

Throughout the course of our partnership with Nurus, our team was featured in a number of news outlets and social media channels. We were proud to share the success of our project and the extent of Nurus' sponsorship with the wider public. Some of the highlights include:

\begin{itemize}
\item An article on our project and partnership with Nurus was featured in the university's newsletter, which reaches thousands of alumni and supporters.
\item Our team was interviewed by a local news outlet, which aired a segment on our project and partnership with Nurus.
\item We shared our progress and results on social media, including posts on Twitter and Instagram that featured Nurus' logo and tagline.
\item We created a promotional video that showcased our partnership with Nurus and highlighted the role that their support played in our success. This video was shared on social media and at a final presentation event.
\end{itemize}

Overall, our media coverage and public relations efforts helped to showcase the extent of Nurus' sponsorship and the success that we achieved together. We are grateful to Nurus for their support and for the opportunity to share our project with a wider audience.

\section{Introduction to Betelgeuse UAV}

The Betelgeuse UAV has been designed and developed by a team of skilled university students who have utilized cutting-edge technology and expertise to create a high-performance unmanned aircraft. With a wingspan of 270cm and a maximum takeoff weight of nearly 12KG, it is a powerful and agile aircraft that can maneuver quickly and efficiently.

Equipped with two T-motor 2826 4S motors and the Pixhawk Cube Orange flight controller, the Betelgeuse UAV is a reliable and precise aircraft that can fly at a maximum speed of nearly 200 kmph and has a maximum flight time of nearly 1 hour. The Jetson Xavier NX companion computer and the Logitech Full HD Webcam allow for accurate detection and tracking of other UAVs in the air.

The Betelgeuse UAV's fully carbon body makes it lightweight and durable, allowing it to fly in most weather conditions. It has a range of up to 10 KM with 900 MHz telemetry and long-range WiFi for communication, making it suitable for a variety of applications.

The primary objective of the Betelgeuse UAV is to track and follow other UAVs for a specific time period of 4 seconds. Its high-performance camera and advanced image processing capabilities enable it to track other UAVs accurately and with high precision. This makes it an ideal choice for applications such as aerial photography, surveillance, and search and rescue operations.

Overall, the Betelgeuse UAV is a remarkable achievement of a university student team, and its advanced technology and precise performance make it an excellent choice for a range of applications.

\section*{Special Thanks}

I would like to express my deepest gratitude to the following people who have supported us throughout this project:

\begin{itemize}
  \item [Name 1] for their guidance and expertise in [specific area].
  \item [Name 2] for their endless encouragement and support.
  \item [Name 3] for their invaluable feedback and constructive criticism.
\end{itemize}

\end{document}
